%% Supplementary material for the Knowledge-Based Systems submission.
%% Compile with pdflatex; uses the same macro files as manuscript.tex.
\documentclass[11pt]{article}
\usepackage[a4paper,margin=2.5cm]{geometry}
\usepackage{amsmath,amssymb}
\usepackage{booktabs}
\usepackage{float}
\usepackage{xurl}
\usepackage[hidelinks]{hyperref}
\renewcommand{\thetable}{S\arabic{table}}

\newcommand{\hlcl}{\textup{\textsc{HLCL}}}
\newcommand{\fdr}{\mathrm{mFDR}}

\input{numbers}
\input{itemcert}
\input{holm}
\input{fdpratio}
\input{perbaseline}
\input{supp}

\title{Supplementary material for\\``Certified label-level selective prediction for hierarchical outputs, with an application to LLM annotation''}
\author{Hongsheng Wang and Fei Wu}
\date{}

\begin{document}
\maketitle

These tables complement Sections~3 and~6 of the main text, and the last two
sections give implementation notes for Section~5 and the annotation prompt of
Section~5.1; table and section numbers without
the prefix S refer to the main text. All tables in this
supplement except Table~\ref{tab:popcheck} use the splits of
Table~3 of the main text. Tables~\ref{tab:depth}, \ref{tab:popcheck}
and~\ref{tab:hypac}--\ref{tab:annotshift} come from a rerun of the deployed rules
from the cached responses on a second machine, which can differ
from the released runs in single splits (see the implementation notes below).

\begin{table}[H]
\centering
\footnotesize
\caption{Error and power of \hlcl{} by depth, mean over 100 splits. $q_d/\alpha$
is the error among the labels accepted at depth $d$ divided by $\alpha$ (--: no
acceptance at that depth); power is the share of correct labels of each depth
that are accepted; stored is the error of all stored labels when a reviewer keeps
the accepted prefix and completes the path, divided by $\alpha$. $q_d$ is
averaged over the splits with at least one acceptance at depth $d$ and is noisy
where power is near zero.}
\label{tab:depth}
\input{depth_table}
\end{table}

\begin{table}[H]
\centering
\footnotesize
\caption{Direct check of the guarantee. The 3,000 items of each setting are
treated as the population; in each of 100 draws, fit and calibration samples are
drawn from them with replacement, and the exact mFDR of the deployed rule is
computed on all 3,000 items. The table gives the share of draws in which it
exceeds $\alpha$; for \hlcl{}, Item-cert.\ (automatic test choice) and
Per-field, the guarantee requires at most $\delta=10\%$.}
\label{tab:popcheck}
\input{popcheck_table}
\end{table}

\begin{table}[H]
\centering
\footnotesize
\setlength{\tabcolsep}{3pt}
\caption{Mean realized label-level error divided by $\alpha$ over 100 splits.
All methods are compared at the same nominal $\alpha$, but they use the error
budget to different degrees.}
\label{tab:realized}
\input{realized_table}
\end{table}

\begin{table}[H]
\centering
\footnotesize
\caption{Paired differences between the tied variant, which decides per label
with one threshold shared by all depths, and Item-cert.$^\ast$, with
Nadeau--Bengio-corrected 95\% intervals over 100 splits. Bold marks significant
differences in favor of the tied variant.}
\label{tab:tied}
\input{tied_table}
\end{table}

\begin{table}[H]
\centering
\footnotesize
\setlength{\tabcolsep}{4pt}
\caption{HyPAC adapted to one model and a person, against \hlcl{}, over 100
splits. viol.: share of splits whose label-level error exceeded $\alpha$
($^\dagger$: more than $\delta=10\%$); error$/\alpha$: wrong over accepted labels,
summed over the splits, divided by $\alpha$; the next two columns give the mean
effort; the last column gives the paired effort difference with
Nadeau--Bengio-corrected 95\% intervals.}
\label{tab:hypac}
\input{hypac_table}
\end{table}

\begin{table}[H]
\centering
\footnotesize
\caption{Effort difference \hlcl{} $-$ Item-cert.$^\ast$ under the default cost
model, under fine-heavy weights $\omega_d\propto d$ and under $c_0=2$, for the
same deployed rules, with Nadeau--Bengio-corrected 95\% intervals. Bold marks
significant differences in favor of \hlcl{}.}
\label{tab:costsens}
\input{costsens_table}
\end{table}

\begin{table}[H]
\centering
\footnotesize
\caption{\hlcl{} under label shift, with the deployed rules of the 100 splits
and no recalibration. The first two columns are for the regular test split; the
last two are for its items in the level-1 classes whose consensus accuracy on the
fit split is below the median. Columns as in Table~\ref{tab:hypac}.}
\label{tab:shift}
\input{shift_table}
\end{table}

\begin{table}[H]
\centering
\footnotesize
\caption{\hlcl{} under a change of annotator on the Blurb collection: the rule
fitted and calibrated on the outputs of the first annotator is applied to the
second annotator's outputs for the same test items, without recalibration.
Columns as in Table~\ref{tab:hypac}.}
\label{tab:annotshift}
\input{annotshift_table}
\end{table}

\section*{Implementation notes}

\paragraph{Label matching (Section~5.1 of the main text).} Labels are compared
after mapping typographic apostrophes to plain ones (annotator B writes
\emph{Children\textquotesingle s Books}, with a straight apostrophe, where the
taxonomy has \emph{Children\textquoteright s Books}).

\paragraph{HyPAC adaptation (Section~5.2 of the main text).} HyPAC selects the
cheapest threshold whose upper confidence bound is at most the target, relying
on the monotonicity of the risk. We instead test the 64 item-score candidates
of Item-cert.\ in a fixed sequence from conservative to permissive, which does
not need that monotonicity, with a one-sided betting $p$-value (HyPAC reports a
betting bound as its best finite-sample variant), and deploy the most
permissive, hence cheapest, certified threshold.

\paragraph{Development of the method (Section~5.4 of the main text).} The
values of $\alpha$ were set per corpus from the consensus accuracy of Table~2
of the main text, which was computed on all 3,000 items. The method and the
grids of $\alpha$ were developed iteratively on the same corpora with
annotators A and B, and the grid of the Blurb collection was changed once, from
$\{0.05,0.1,0.15,0.2\}$ to $\{0.02,0.05,0.1,0.15\}$, after early runs. The
candidate order and the test choice use the in-sample scores of the fit split,
which we preferred to a cross-fitted variant in a pilot; both are functions of
the fit split only and do not affect validity. Annotator C was added after the
code was frozen, so its four settings are a check on data not used in
development.

\paragraph{Reruns on other machines (Section~5.4 of the main text).} Rerunning
the experiments on a different machine can change single splits through
floating-point ties among candidate $p$-values; in our reruns of all
\stat{s@nsettings} settings, 48 of 2,400 splits of \hlcl{} changed, mean effort
changed by at most $0.008$ (DBPedia at $\alpha=0.02$) and by at most $0.002$
elsewhere, and the significance counts of Section~6.2 of the main text were
unchanged.

\section*{Annotation prompt}

Every item was annotated with the following prompt, sent as a single user
message with temperature 1.0, at most 160 output tokens and JSON output mode
where the API supports it. \texttt{\{taxonomy\}} lists the taxonomy with one line per deepest parent, in
the form ``\texttt{a > b > leaf1 | leaf2 | ...}''; \texttt{\{text\}} is the item
text truncated to 1,500 characters. The prompt is shown for a three-level
taxonomy; for two and four levels, the header and the JSON keys have two or four
levels. The full taxonomy is included in every prompt.

\begin{footnotesize}
\begin{verbatim}
Classify the text into exactly one path of the taxonomy below
(coarse to fine). Use label names exactly as written.

TAXONOMY (each line: level-1 > level-2 > level-3 options):
{taxonomy}

TEXT:
{text}

Answer with JSON only: {"l1", "l2", "l3", "c1", "c2", "c3"} where c* is
your probability (0-1) that the label at that level is correct.
\end{verbatim}
\end{footnotesize}

\end{document}
